% GENERATED FILE. Edit canonical lessons, then run npm run generate:books.
% canonical-source-hash: fnv1a64:8192ac2a35012385
% canonical-lessons: ZH-C448-tanwei, ZH-C448-shichang, ZH-W448-shi-market, ZH-C448-shiqu, ZH-C448-zaoshi, ZH-C448-baitan

\chapter{At the Market}
\label{ch:zh-at-the-market}
\hlchaptermodality{448}

\begin{chapteropening}
\textbf{By the end of this chapter:} \emph{I can talk about stalls, find the market, stay in town, go to a morning market and run a stall.}

Complete the last lesson of chapter 448: I can talk about stalls, find the market, stay in town, go to a morning market and run a stall.
\end{chapteropening}

\section[tānwèi]{tānwèi — market stall}

\label{lesson:ZH-C448-tanwei}



\begin{quote}
Before the new one: say the Mandarin for to take the long way round, then the Mandarin for the pavement beside the road.
\end{quote}

\subsection*{You'll want to know: tānwèi}
\emph{tānwèi} — \textquotedblleft{}market stall\textquotedblright{}.

\begin{quote}
\emph{zhège tānwèi mài shuǐguǒ} — This stall sells fruit.
\end{quote}

Its characters come later; for now it is a word for the ear.

\begin{cousinweb}
\emph{tānwèi} is built from \emph{tān} 'to spread out' and \emph{wèi} 'a place', both inherited from Middle Chinese, \emph{than} and \emph{hwijH} (Guangyun, in Baxter's transcription), which glosses the first 'to open out' and the second as a place in a row.
\end{cousinweb}

\begin{grammarlens}[title={What you've built}]
One more word for getting around.
\end{grammarlens}

\subsection*{Your turn}
\begin{itemize}
\raggedright
  \item \emph{Say it:} \emph{tānwèi}
  \item \emph{Say it:} \emph{tānwèi}, once more
  \item \emph{Say it:} \emph{zhège tānwèi mài shuǐguǒ} — This stall sells fruit.
  \item \emph{Recall:} say the Mandarin for to take the long way round, then the Mandarin for the pavement beside the road, then say \emph{tānwèi} again
\end{itemize}

\begin{tcolorbox}[breakable,colback=teal!4,colframe=teal!35!black,title={Before you move on}]
What does tānwèi mean? (\textquotedblleft{}Market stall\textquotedblright{}.) Say it once more.
\end{tcolorbox}
\section[\texorpdfstring{\zh{市场} (shìchǎng)}{市场 (shìchǎng)}]{\zh{市场} (shìchǎng) — market; marketplace}

\label{lesson:ZH-C448-shichang}



\begin{quote}
Before the new one: say the Mandarin for the pavement beside the road, then the Mandarin for a market stall.
\end{quote}

\subsection*{You'll want to know: \zh{市场}}
\textbf{\zh{市场}} — \emph{shìchǎng} — \textquotedblleft{}market; marketplace\textquotedblright{}.

\begin{quote}
\textbf{\zh{去市场怎么走}} \emph{qù shìchǎng zěnme zǒu} — How do I get to the market?
\end{quote}

The new character, in pieces: \zh{市} — five strokes: a dot, a long bar, then a short vertical, a stroke that turns down and hooks, and a long vertical through the middle. The next lesson puts it in your hand.

\begin{cousinweb}
\textbf{\zh{市场}} (\emph{shìchǎng}) is built from \zh{市} 'a market' and \zh{场} 'an open ground', both inherited from Middle Chinese, \emph{dzyiX} and \emph{drjang} (Guangyun, in Baxter's transcription), which glosses the first, after an older dictionary, as where people go to buy and sell, and the second as a threshing floor.
\end{cousinweb}

\begin{grammarlens}[title={What you've built}]
One more word for getting around.
\end{grammarlens}

\subsection*{Your turn}
\begin{itemize}
\raggedright
  \item \emph{Say it:} \emph{shìchǎng}
  \item \emph{Say it:} \emph{shìchǎng}, once more
  \item \emph{Say it:} \emph{qù shìchǎng zěnme zǒu} — How do I get to the market?
  \item \emph{Recall:} say the Mandarin for the pavement beside the road, then the Mandarin for a market stall, then say \emph{shìchǎng} again
\end{itemize}

\begin{tcolorbox}[breakable,colback=teal!4,colframe=teal!35!black,title={Before you move on}]
What does \zh{市场} mean? (\textquotedblleft{}Market; marketplace\textquotedblright{}.) Say it once more.
\end{tcolorbox}
\section[\texorpdfstring{\zh{市} (shì)}{市 (shì)}]{\zh{市} — a dot and a bar over three strokes}

\label{lesson:ZH-W448-shi-market}



\begin{quote}
The word first, then the character inside it.

\begin{itemize}
\raggedright
  \item \emph{Recall:} say \emph{shìchǎng} — \textbf{R1}, one lesson back, before it decays
\end{itemize}

\begin{quote}
\zh{场} \zh{方}
\end{quote}

\emph{Read it:} each one aloud, with nothing said first — \emph{chǎng}, \emph{fāng}
\end{quote}

\subsection*{Script — a dot and a bar over three strokes}
\hlblockfigure{\detokenize{figures/ZH-W448-shi-market-filmstrip.pdf}}{How \zh{市} is written, stroke by stroke}

\begin{quote}
\zh{市}
\end{quote}

\textbf{One piece, five strokes, top to bottom.}

\textbf{A dot on top, then a long bar}, as at the top of \zh{方}.

\textbf{Under the bar, three strokes like a small hanging cloth.} A short vertical on the left. A stroke that runs right and turns down, hooking to the left at the foot. Then a long vertical through the middle, from the bar down past the hook.

2 + 3 = five.

\subsection*{Your turn}
\begin{itemize}
\raggedright
  \item \emph{Write it:} \zh{市}, all five, with the model in front of you
  \item \emph{Cover:} write it again from memory, and count the pieces, not the strokes
\end{itemize}

\begin{tcolorbox}[breakable,colback=teal!4,colframe=teal!35!black,title={Before you move on}]
How many strokes in \zh{市}? (\textbf{Five}.) What comes first? (\textbf{A dot}.) What is the last stroke? (\textbf{The long vertical through the middle}.)

Source: \href{https://raw.githubusercontent.com/chanind/hanzi-writer-data/68d10a4b21150cae5e1ebbd223eed289cf32d90c/data/\%E5\%B8\%82.json}{Hanzi Writer Data \zh{市}, pinned stroke paths and medians}.
\end{tcolorbox}
\section[\texorpdfstring{\zh{市区} (shìqū)}{市区 (shìqū)}]{\zh{市区} (shìqū) — the built-up part of a city; downtown}

\label{lesson:ZH-C448-shiqu}



\begin{quote}
Before the new one: say the Mandarin for a market stall, then the Mandarin for a marketplace.
\end{quote}

\subsection*{You'll want to know: \zh{市区}}
\textbf{\zh{市区}} — \emph{shìqū} — \textquotedblleft{}the built-up part of a city; downtown\textquotedblright{}.

\begin{quote}
\emph{wǒmen zhù zài shìqū de yì jiā jiǔdiàn} — We're staying at a hotel in town.
\end{quote}

\begin{cousinweb}
\textbf{\zh{市区}} (\emph{shìqū}) is built from \zh{市} 'a market town' and \zh{区} 'an area', both inherited from Middle Chinese, \emph{dzyiX} and \emph{khju} (Guangyun, in Baxter's transcription).
\end{cousinweb}

\begin{grammarlens}[title={What you've built}]
One more word for getting around.
\end{grammarlens}

\subsection*{Your turn}
\begin{itemize}
\raggedright
  \item \emph{Say it:} \emph{shìqū}
  \item \emph{Say it:} \emph{shìqū}, once more
  \item \emph{Say it:} \emph{wǒmen zhù zài shìqū de yì jiā jiǔdiàn} — We're staying at a hotel in town.
  \item \emph{Recall:} say the Mandarin for a market stall, then the Mandarin for a marketplace, then say \emph{shìqū} again
\end{itemize}

\begin{tcolorbox}[breakable,colback=teal!4,colframe=teal!35!black,title={Before you move on}]
What does \zh{市区} mean? (\textquotedblleft{}The built-up part of a city; downtown\textquotedblright{}.) Say it once more.
\end{tcolorbox}
\section[\texorpdfstring{\zh{早市} (zǎoshì)}{早市 (zǎoshì)}]{\zh{早市} (zǎoshì) — morning market; a market held at dawn}

\label{lesson:ZH-C448-zaoshi}



\begin{quote}
Before the new one: say the Mandarin for a marketplace, then the Mandarin for the built-up part of a city.
\end{quote}

\subsection*{You'll want to know: \zh{早市}}
\textbf{\zh{早市}} — \emph{zǎoshì} — \textquotedblleft{}morning market; a market held at dawn\textquotedblright{}.

\begin{quote}
\textbf{\zh{我们早上去早市买菜}} \emph{wǒmen zǎoshang qù zǎoshì mǎi cài} — In the morning we go to the early market to buy vegetables.
\end{quote}

\begin{cousinweb}
\textbf{\zh{早市}} (\emph{zǎoshì}) is built from \zh{早} 'early' and \zh{市} 'a market', both inherited from Middle Chinese, \emph{tsawX} and \emph{dzyiX} (Guangyun, in Baxter's transcription), which glosses the first 'the morning'.
\end{cousinweb}

\begin{grammarlens}[title={What you've built}]
One more word for getting around.
\end{grammarlens}

\subsection*{Your turn}
\begin{itemize}
\raggedright
  \item \emph{Say it:} \emph{zǎoshì}
  \item \emph{Say it:} \emph{zǎoshì}, once more
  \item \emph{Say it:} \emph{wǒmen zǎoshang qù zǎoshì mǎi cài} — In the morning we go to the early market to buy vegetables.
  \item \emph{Recall:} say the Mandarin for a marketplace, then the Mandarin for the built-up part of a city, then say \emph{zǎoshì} again
\end{itemize}

\begin{tcolorbox}[breakable,colback=teal!4,colframe=teal!35!black,title={Before you move on}]
What does \zh{早市} mean? (\textquotedblleft{}Morning market; a market held at dawn\textquotedblright{}.) Say it once more.
\end{tcolorbox}
\section[bǎitān]{bǎitān — to set up a street stall}

\label{lesson:ZH-C448-baitan}



\begin{quote}
Before the new one: say the Mandarin for the built-up part of a city, then the Mandarin for a market held at dawn.
\end{quote}

\subsection*{You'll want to know: bǎitān}
\emph{bǎitān} — \textquotedblleft{}to set up a street stall\textquotedblright{}.

\begin{quote}
\emph{tā měitiān zài shìchǎng bǎitān} — He runs a stall at the market every day.
\end{quote}

Its characters come later; for now it is a word for the ear.

\begin{cousinweb}
\emph{bǎitān} is built from \emph{bǎi} 'to set out' and \emph{tān} 'to spread out', both inherited from Middle Chinese, \emph{peaX} and \emph{than} (Guangyun, in Baxter's transcription), which glosses the first 'to push aside' and the second 'to open out'; the sense 'to set out goods' for the first is not given there.
\end{cousinweb}

\begin{grammarlens}[title={What you've built}]
One more word for getting around.
\end{grammarlens}

\subsection*{Your turn}
\begin{itemize}
\raggedright
  \item \emph{Say it:} \emph{bǎitān}
  \item \emph{Say it:} \emph{bǎitān}, once more
  \item \emph{Say it:} \emph{tā měitiān zài shìchǎng bǎitān} — He runs a stall at the market every day.
  \item \emph{Recall:} say the Mandarin for the built-up part of a city, then the Mandarin for a market held at dawn, then say \emph{bǎitān} again
\end{itemize}

\begin{tcolorbox}[breakable,colback=teal!4,colframe=teal!35!black,title={Before you move on}]
What does bǎitān mean? (\textquotedblleft{}To set up a street stall\textquotedblright{}.) Say it once more.
\end{tcolorbox}
